\documentclass[UTF8,a4paper,12pt]{ctexart}

\usepackage{amsmath,amssymb,bm}
\usepackage{booktabs}
\usepackage{geometry}
\usepackage{longtable}
\usepackage{array}
\usepackage{enumitem}
\usepackage{microtype}
\usepackage{hyperref}

\geometry{left=2.6cm,right=2.6cm,top=2.5cm,bottom=2.5cm}
\setlength{\parindent}{2em}
\setlength{\parskip}{0.35em}
\renewcommand{\arraystretch}{1.35}
\setlist{nosep}
\hypersetup{hidelinks}

\title{方法与实验部分}
\author{}
\date{}

\begin{document}

\maketitle

\section{3 提出的模型}

本研究面向大型多线地铁换乘站的整体应急疏散，构建由站内设施网络、多来源客流到达过程、共享设施服务队列和动态路径组织共同组成的疏散模型。模型的出发点不是比较某一条算法路径是否更短，而是刻画多线路站台客流、列车到达客流和换乘客流在站内共享楼扶梯、换乘通道、闸机和出口处的汇聚过程，并在此基础上组织能够被实际容量执行的疏散路径。参考论文在模型部分依次建立消防路线鲁棒优化模型、节点通行时间预测模型和综合评价方法；对应到本文，模型部分依次建立多线换乘站疏散网络、设施到达负荷预测与动态路径组织模型，以及整体疏散效果评价方法。

\subsection{3.1 多线换乘站疏散网络模型}

\subsubsection{3.1.1 站内设施网络}

将龙阳路站公共区抽象为有向图
\[
G=(V,E),
\]
其中，节点集合 \(V\) 表示线路站台、站厅、换乘连接节点、楼梯、扶梯、闸机、出口前区域和地面出口，边集合 \(E\) 表示相邻空间或设施之间允许乘客移动的方向性连接。每条边 \(e=(i,j)\) 包含长度 \(l_e\)、有效宽度 \(w_e\)、行进方向、设施类型和容量属性；每个节点 \(v\) 包含面积 \(A_v\)、当前人数 \(N_v(t)\)、节点类型和空间接收上限。

大型换乘站的关键特征在于，同一项设施往往不是服务单一线路，而是同时承接多个来源方向。例如，站台通向站厅的楼扶梯会同时接收站台出站客流和换乘客流，换乘通道会接收来自不同线路的交叉流，闸机和出口前区域则可能成为多个线路客流的共同终端。因此，模型不把每条边的容量独立计算，而是建立边到共享设施资源的映射。若边 \(e\) 需要消耗共享设施 \(r\) 的通行能力，则记为 \(\rho(e)=r\)；若边只受普通空间移动约束，则记为 \(\rho(e)=0\)。映射到同一设施 \(r\) 的多条边统一使用该设施的服务能力。

\subsubsection{3.1.2 多来源客流与换乘客流}

设客流来源组集合为 \(\mathcal{G}\)。来源组用于区分不同线路站台、列车到达、站厅滞留和跨线换乘客流。对换乘站而言，换乘客流不是独立于整体疏散之外的评价对象，而是决定站内设施负荷分布的重要来源。其特殊性在于：换乘客流的起点、方向和可选出口可能跨越多个线路公共区，并且在疏散过程中会与站台出站客流和列车到达客流共同占用楼扶梯、换乘通道和闸机等设施。

客流以自然到达批次表示。一个批次记为
\[
b=(g,n_b,x_b,t_b),
\]
其中，\(g\) 为来源组，\(n_b\) 为批次人数，\(x_b\) 为当前节点，\(t_b\) 为批次进入路径决策的时间。采用批次而不是单个乘客作为路径决策单元，可以保留列车到达、站台释放和换乘流进入公共区的时间结构，也能使后续路径搜索显式记录不同批次对同一设施的到达先后关系。

\subsubsection{3.1.3 共享设施服务队列}

对具有服务能力限制的设施 \(r\)，设其服务率为 \(\mu_r\)，单位为人每秒；仿真步长为 \(\Delta t\)。在时间步 \(t\) 内，设施 \(r\) 可释放的整数服务能力为
\[
C_r(t)=\left\lfloor \mu_r\Delta t+\xi_r(t)\right\rfloor ,
\]
其中，\(\xi_r(t)\) 为上一时间步未形成完整整数人的容量余量。设 \(A_r(t)\) 为时间步 \(t\) 内请求进入设施 \(r\) 的人数，\(Q_r(t)\) 为该设施入口等待队列，则队列递推为
\[
Q_r(t+\Delta t)=\mathrm{max}\left(Q_r(t)+A_r(t)-C_r(t),0\right).
\]

这一表达式的意义在于，楼梯、扶梯、通道和闸机不是路径图上的普通连接线，而是具有服务释放过程的瓶颈资源。若多个来源组在同一时间步到达同一资源，它们应共同竞争同一 \(C_r(t)\)，未获得服务能力的人数应停留在上游等待，而不能因为路径搜索已经给出方向就被计为通过。

\subsubsection{3.1.4 空间接收与行走时间}

设施服务能力还受到下游空间接收条件限制。设节点 \(v\) 的空间接收上限为 \(\bar N_v\)，当前占用人数为 \(N_v(t)\)，当前时间步已经被其他移动请求预留进入的人数为 \(B_v(t)\)，则新进入人数 \(y_v(t)\) 需要满足
\[
N_v(t)+B_v(t)+y_v(t)\leq \bar N_v .
\]
当下游节点接收不足时，即使上游设施仍有名义服务能力，实际移动也可能被阻滞。这一约束用于描述高负荷换乘站中常见的回溢现象，即队列不只在闸机或楼梯入口形成，也可能从下游通道、站厅或出口前区域向上游扩散。

对于普通移动边，行走时间由边长和速度决定。当前实现采用速度--密度关系计算边内通行时间：
\[
v(k)=
\begin{cases}
v_f, & k\leq k_f,\\
\max(v_f-ak,0), & k_f<k<k_j,\\
0, & k\geq k_j .
\end{cases}
\]
其中，\(k\) 为边内客流密度，\(v_f\) 为自由流速度，\(k_f\) 为自由流密度阈值，\(a\) 为速度下降斜率，\(k_j\) 为拥挤密度上限。本文采用的 \(v_f=1.427\ \mathrm{m/s}\)、\(k_f=0.2\ \mathrm{p/m^2}\)、\(a=0.3549\) 和 \(k_j=4.0\ \mathrm{p/m^2}\) 来自对改进 A* 疏散仿真文献中速度--密度关系的复现，不作为龙阳路站实测速度。该关系用于计算物理行走时间，但本文的实验边界分析不以速度--密度阈值为核心变量。

\subsection{3.2 面向换乘客流汇聚的到达负荷预测与路径组织模型}

\subsubsection{3.2.1 设施到达负荷预测}

在低负荷条件下，当前队列与到达时队列之间的差异可能较小；在高负荷换乘站中，多个来源批次会在不同时间先后到达同一共享设施，当前时刻尚未出现的排队可能在候选批次到达时已经形成。因此，路径评价应预测批次到达关键设施时的服务状态，而不是只读取当前时刻的节点人数。

对共享设施 \(r\)，定义未来确定到达事件集合
\[
\mathcal{A}_r(t)=\{(\eta_j,n_j)\},
\]
其中，\(\eta_j\) 为已确定批次预计进入设施 \(r\) 的时刻，\(n_j\) 为对应人数。对候选批次 \(b\)，若其预计在时刻 \(\eta\) 到达设施 \(r\)，则设施到达队列估计为
\[
\widehat Q_r(\eta)=\mathcal{Q}\left(Q_r(t),\mathcal{A}_r(t),\mu_r,t,\eta\right).
\]
队列演化算子 \(\mathcal{Q}\) 在 \(t\) 到 \(\eta\) 的时间区间内释放服务能力，并在已登记到达事件发生时加入相应人数。若 \(\mu_r>0\)，候选批次在设施 \(r\) 处的预计等待时间为
\[
W_r(\eta)=\frac{\widehat Q_r(\eta)}{\mu_r}.
\]

该预测不是机器学习模型，而是基于批次到达事件和服务能力的确定性递推。它在本文中的作用相当于参考论文中“节点通行时间预测模型”的位置：为上层路径组织提供关键设施通行时间和等待时间依据。

\subsubsection{3.2.2 时间依赖路径代价}

对批次 \(b\) 的候选路径 \(P=(e_1,e_2,\ldots,e_m)\)，设第 \(i\) 条边的预计进入时刻为 \(\eta_i\)，对应共享设施为 \(r_i=\rho(e_i)\)。路径代价定义为
\[
J_b(P)=\sum_{i=1}^{m}\left[\tau_{e_i}(\eta_i)+W_{r_i}(\eta_i)+S_{r_i}(n_b)+H_{e_i}(\eta_i)\right],
\]
其中，\(\tau_{e_i}\) 为物理行走时间，\(W_{r_i}\) 为到达设施时的预计排队等待，\(S_{r_i}(n_b)\) 为当前批次自身占用服务能力形成的平均服务时间，\(H_{e_i}\) 为下游空间接收不足或回溢造成的预计等待。若边不对应共享设施，则相关服务等待项取 0。

路径代价按预计到达时刻递推。也就是说，前一段路径上的行走和等待会改变后一段路径的到达时间，后一段设施的等待又取决于该新的到达时间。由此，本文的路径组织不是静态最短路，也不是只根据当前密度重新计算边权，而是面向设施到达负荷的时间依赖路径搜索。

\subsubsection{3.2.3 路径重评估约束}

动态组织不等于乘客在移动途中任意改变方向。本文只对尚未进入在途执行状态、仍处于明确选择阶段的批次进行路径重评估。已进入边内移动或已经获得设施入口服务能力的批次继续执行已承诺的移动。

设保留当前路径的剩余代价为 \(C_{\mathrm{stay}}\)，候选切换路径代价为 \(C_{\mathrm{switch}}\)，相对收益为
\[
R=\frac{C_{\mathrm{stay}}-C_{\mathrm{switch}}}{C_{\mathrm{stay}}}.
\]
仅当候选路径有效、下一步移动方向发生变化、批次仍处于可重评估阶段、\(C_{\mathrm{switch}}<C_{\mathrm{stay}}\)，且 \(R\) 超过防摆动阈值时，路径切换才被接受。该阈值是算法稳定性设置，不是行人行为参数。

\subsection{3.3 车站疏散效果综合评价模型}

评价指标服务于整体疏散，而不是单独证明某一个算法优于另一个算法。本文从四个层面评价路径组织效果。

第一，整体疏散效率。采用 \(T_{95}\)、\(T_{100}\)、平均疏散时间、累计停滞人秒和有效疏散速度描述全站清空过程。其中，\(T_{95}\) 用于刻画系统进入尾部清空阶段前的效率，\(T_{100}\) 用于描述最后一批乘客离站时间，累计停滞人秒用于衡量乘客在排队、阻塞和等待状态中的暴露程度。

第二，关键设施瓶颈。对楼扶梯、换乘通道、闸机和出口前区域记录通过人数、峰值队列、队列人秒、首次排队时间、最后排队时间和利用率。若某一方法降低系统尾部时间但增加局部设施队列，该变化必须在该层指标中被如实呈现。

第三，来源组与换乘客流分解。按线路站台、列车到达、站厅和换乘来源统计出口分配、设施通过量和清空时间，用于解释不同来源客流在哪些共享设施上汇聚，以及换乘客流对全站尾部疏散的贡献。

第四，跨平台一致性。将宏观网络模型的路径组织结果导入 Pathfinder 微观仿真，比较疏散曲线、拥堵空间分布、出口利用和乘客拥堵时间。Pathfinder 不是现场观测真值，而是用于检验宏观路径组织结果在微观行人运动尺度上是否呈现相近趋势。

\section{4 求解方法}

本节说明如何将第 3 节的到达负荷预测和路径组织模型转化为可执行算法。参考论文第 4 节通过模型转换和求解算法承接鲁棒优化模型；本文对应地将时间依赖队列评价转换为批次级路径搜索，并通过共享容量执行层保证路径决策与设施服务过程一致。

\subsection{4.1 时间依赖路径搜索状态转换}

对批次 \(b\)，路径搜索标签定义为
\[
L=(v,\eta,g,P,\pi),
\]
其中，\(v\) 为当前节点，\(\eta\) 为预计到达当前节点的时刻，\(g\) 为累计代价，\(P\) 为已经形成的局部路径，\(\pi\) 为前驱信息。标签从节点 \(u\) 沿边 \(e=(u,v)\) 扩展时，先计算物理行走时间，再根据预计进入时刻查询共享设施队列和下游空间接收状态，得到新的到达时刻与累计代价。

与静态 A* 相比，本文的标签不能只用节点编号合并。两个标签即使到达同一节点，只要预计到达时刻不同，后续设施的队列状态就可能不同，因此需要保留具有不同到达时刻和不同资源使用后果的标签。静态自由流时间仅作为启发式下界使用，不替代实际路径代价。

\subsection{4.2 队列感知动态路径组织算法}

算法以离散时间步推进，每一时间步包含状态更新、批次决策、资源执行和指标记录四个环节。

\begin{enumerate}
    \item 更新上一时间步已经完成移动的乘客，释放其进入目标节点或设施队列。
    \item 根据站台释放、列车到达和换乘来源规则生成本时间步进入决策阶段的自然到达批次。
    \item 按到达时间、来源组、当前节点和批次编号对可决策批次排序。
    \item 对每个批次执行时间依赖 A* 搜索，计算候选路径的预计行走时间、设施等待时间和空间阻塞等待。
    \item 将已确定路径的整数批次登记为未来设施到达事件，使后续批次能够看到本轮已经承诺的设施负荷。
    \item 汇总所有上游节点对同一共享设施的移动请求，并按照该设施的整数服务能力统一分配通过人数。
    \item 检查下游节点接收能力和边内空间状态，将未获得容量或未能进入下游的人数保留在上游等待状态。
    \item 记录全站疏散、关键设施队列、来源组位置、出口分配和路径重评估事件。
\end{enumerate}

本文提出的方法记为 AdaptiveQueueAwareAStar，即队列感知动态路径组织方法。它的核心不是简单增加 A* 的计算次数，而是在搜索过程中使用候选批次到达共享设施时的预计队列状态。

\subsection{4.3 计算对照方法与共同执行层}

ImprovedAStar 在本文中仅作为计算对照方法。它在同一车站网络、同一客流输入和同一速度--密度关系下计算路径，但不使用未来设施到达事件集合，不预测候选批次到达设施时的排队状态，也不执行基于收益阈值的路径重评估。ImprovedAStar 不代表龙阳路站运营方制定或发布的疏散预案。

为保证比较的可解释性，ImprovedAStar 与本文方法共用同一共享容量执行层。路径搜索只给出候选移动方向，实际通过人数仍由楼扶梯、通道、闸机和出口的服务能力以及下游空间接收条件决定。由此，两种方法的差异被限定在路径组织逻辑，而不是设施容量、客流输入或物理执行规则。

\section{5 案例研究与实验设计}

本文以龙阳路站为案例进行实验验证。公开资料显示，龙阳路站承担 2 号线、7 号线、16 号线、18 号线及磁浮线之间的换乘功能，车站空间层级和换乘组织较为复杂。模型根据 CAD 图纸和站内拓扑配置建立站台、站厅、换乘通道、楼扶梯、闸机和出口网络，设施宽度、闸机数量和出口宽度在 CAD 核对完成后进入最终参数表。本节实验设计按照参考论文案例研究的层次展开：先验证中间预测模块，再比较模型方案性能，最后将路径组织结果放回完整疏散仿真和 Pathfinder 微观模型中进行综合评价。

\subsection{5.1 设施到达负荷预测结果分析}

本节对应参考论文对“节点通行时间预测结果”的分析。由于本文不采用训练型节点时间预测器，验证对象改为基于批次到达事件的设施到达负荷预测。该实验回答三个问题：第一，模型是否能够把前序批次已经承诺的到达事件计入后续批次的设施等待；第二，预测等待是否会改变候选路径排序；第三，预测得到的关键设施压力是否能在执行层的实际队列中得到对应体现。

实验选取换乘通道、楼扶梯、闸机和出口前区域中的关键共享设施作为分析对象。对每个设施记录实际队列、未来确定到达人数、预测等待时间、实际通过人数和未获得服务能力人数。对代表性来源组，进一步分解其候选路径代价，包括物理行走时间、设施等待时间、空间回溢等待和总路径代价。结果冻结后，本节将给出关键设施队列曲线和代表性批次路径代价分解图，用于说明队列预测如何改变路径组织判断。

\begin{table}[htbp]
\centering
\caption{设施到达负荷预测分析指标}
\small
\begin{tabular}{p{0.20\textwidth}p{0.34\textwidth}p{0.34\textwidth}}
\toprule
分析对象 & 指标 & 作用 \\
\midrule
关键楼扶梯与换乘通道 & 预测到达人数、实际队列、峰值等待、最后排队时间 & 判断多来源客流是否在垂向设施和换乘连接处集中汇聚 \\
闸机与出口前区域 & 通过人数、队列人秒、出口分配、来源组构成 & 判断出口方向选择是否造成局部服务压力 \\
代表性来源组批次 & 路径代价分解、候选路径排序、是否触发重评估 & 解释路径改变来自行走距离、设施等待还是空间回溢 \\
\bottomrule
\end{tabular}
\end{table}

\subsection{5.2 疏散路径组织方案性能分析}

\subsubsection{5.2.1 实验场景}

本文设置低负荷和高负荷两类疏散场景。低负荷常规应急场景包含各线路站台等待、站厅和换乘基础客流，输入总人数为 2187 人。该场景用于检验在设施竞争较弱时，队列感知路径组织是否保持稳定，是否因过度预测引入不必要绕行或局部出口偏置。

高负荷满载列车叠加场景在基础客流上加入多线路列车到达客流，输入总人数为 17905 人。该场景用于放大多线路站台客流、列车客流和换乘客流在共享设施上的汇聚效应，检验到达负荷预测能否降低整体等待暴露和尾部滞留。

两类场景均比较 ImprovedAStar 与 AdaptiveQueueAwareAStar。两种方法使用相同的站内网络、客流输入、设施服务能力、速度--密度关系、空间接收约束和终止规则。最终结果冻结前，正文不写入百分比改善或最终优劣结论。

\subsubsection{5.2.2 整体疏散效率比较}

整体性能首先比较 \(T_{95}\)、\(T_{100}\)、平均疏散时间、累计停滞人秒、平均移动时间、总移动距离和运行时间。低负荷场景重点观察方法是否保持稳定；高负荷场景重点观察尾部清空、停滞等待和计算代价。若某方法降低等待但增加移动距离，或者缩短系统尾部但增加局部设施队列，正文需要同时报告这些权衡。

\begin{table}[htbp]
\centering
\caption{低负荷与高负荷场景下的整体疏散性能记录表}
\small
\begin{tabular}{p{0.16\textwidth}p{0.18\textwidth}p{0.10\textwidth}p{0.12\textwidth}p{0.12\textwidth}p{0.14\textwidth}p{0.12\textwidth}}
\toprule
场景 & 方法 & 人数 & \(T_{95}\) & \(T_{100}\) & 累计停滞人秒 & 总移动距离 \\
\midrule
低负荷常规应急 & ImprovedAStar & 2187 & 待冻结 & 待冻结 & 待冻结 & 待冻结 \\
低负荷常规应急 & AdaptiveQueueAwareAStar & 2187 & 待冻结 & 待冻结 & 待冻结 & 待冻结 \\
高负荷满载列车叠加 & ImprovedAStar & 17905 & 待冻结 & 待冻结 & 待冻结 & 待冻结 \\
高负荷满载列车叠加 & AdaptiveQueueAwareAStar & 17905 & 待冻结 & 待冻结 & 待冻结 & 待冻结 \\
\bottomrule
\end{tabular}
\end{table}

\subsubsection{5.2.3 关键设施与来源组分解}

在整体指标之后，进一步分析关键设施队列和来源组分解。该分析不只是补充图表，而是用于解释整体差异的物理原因。若高负荷场景的尾部清空时间发生变化，需要追踪这种变化是否来自某些换乘通道、楼扶梯或闸机的持续排队减少；若系统尾部时间缩短但某一线路或某一区域排队增加，也需要在该层结果中说明局部代价。

来源组分解重点关注换乘客流与站台、列车来源客流在共享设施上的叠加关系。对每一类来源组统计其主要路径链、关键设施通过量、出口选择和清空时间。通过该分析可以回答：换乘客流是否是某些瓶颈形成的重要组成部分；路径组织改变是否改变了换乘客流到达关键设施的时间分布；以及这种变化是否影响全站整体疏散效率。

\subsection{5.3 路径组织与乘客疏散综合效果分析}

\subsubsection{5.3.1 换乘客流汇聚机制}

综合效果分析首先回到大型换乘站的空间机制。对高负荷场景，按照关键设施队列人秒和峰值队列识别主要瓶颈，再将瓶颈设施的客流来源分解为不同线路站台、列车到达和换乘来源。若同一瓶颈由多个来源共同构成，说明该设施是全站疏散组织中的共享限制环节；若某一方法只是把等待从一个设施转移到另一个设施，则应在路径链和设施队列结果中明确呈现。

\subsubsection{5.3.2 适用边界分析}

本文保留适用边界分析，但不把速度--密度阈值作为主实验变量。速度--密度关系属于行走时间模型参数，对本文核心问题的解释力有限。更能反映多线换乘站疏散组织边界的变量包括三类：换乘客流比例、列车到达重叠程度和关键共享设施服务能力。

换乘客流比例用于检验方法是否主要在换乘客流显著参与设施竞争时发挥作用；列车到达重叠程度用于检验多来源客流同时释放对设施排队和尾部清空的影响；关键共享设施服务能力用于检验方法在瓶颈强弱变化下的适用范围。若某一变量变化后两种方法差异明显缩小，说明该变量对应的场景条件是方法适用性的边界。

\begin{table}[htbp]
\centering
\caption{适用边界分析的变量设置}
\small
\begin{tabular}{p{0.22\textwidth}p{0.28\textwidth}p{0.38\textwidth}}
\toprule
变量类别 & 设置方式 & 分析目的 \\
\midrule
换乘客流比例 & 在总客流规模基本一致的前提下调整换乘来源占比 & 判断换乘客流汇聚是否是方法收益的重要来源 \\
列车到达重叠程度 & 调整多线路列车客流释放的时间间隔 & 判断同步到达和错峰到达对共享设施排队的影响 \\
关键共享设施服务能力 & 对主要楼扶梯、换乘通道或闸机服务率进行扰动 & 判断路径组织效果对瓶颈强度的依赖程度 \\
\bottomrule
\end{tabular}
\end{table}

\subsubsection{5.3.3 Pathfinder 微观仿真对照}

Pathfinder 对照用于检验宏观网络模型的路径组织结果在微观行人仿真中是否呈现一致趋势。本文将宏观模型导出的来源组、路径链和路径人数转换为 Pathfinder 中的行人组和行为路径，并在相同车站几何和出口设置下运行微观仿真。对照指标包括累计疏散曲线、尾部清空时间、拥堵时间、关键设施附近空间拥挤分布和出口利用。

该对照不把 Pathfinder 视为真实车站观测，也不以 Pathfinder 结果替代宏观模型结果。其作用是检查：宏观模型识别出的瓶颈位置是否在微观轨迹中出现相近拥堵；路径组织引起的出口和设施使用变化是否能在微观模型中复现；若两者存在差异，则进一步从导航网格、个体避让、门选择规则和容量映射差异解释。

\subsubsection{5.3.4 结果解释顺序}

结果冻结后，案例研究按照以下顺序写作：首先报告设施到达负荷预测结果，证明模型的中间预测模块参与并改变路径评价；其次比较低负荷和高负荷场景下两种路径组织方法的整体疏散效率；再次通过关键设施和来源组分解解释整体差异的站内机制，重点说明换乘客流与其他来源客流在共享设施上的汇聚关系；随后通过适用边界分析说明方法在不同换乘比例、到达重叠和瓶颈强度下的有效范围；最后用 Pathfinder 微观仿真对宏观疏散趋势进行交叉验证。

\section*{参考文献占位}

\begin{enumerate}
    \item 中国城市轨道交通协会。2026 年上半年中国内地城轨交通线路概况。2026。
    \item 上海市交通委员会。还有 2 天！两条轨交线开通初期运营，上海地铁又添“世界第一”。2021。
    \item 蒙盾，胡卓，张华军。基于改进 A* 算法的多层邮轮疏散系统仿真。系统仿真学报，2022，34(6):1375--1382。
    \item Yang X. 等。Path planning for guided passengers during evacuation in subway station based on multi-objective optimization。Applied Mathematical Modelling，2022。
    \item Zhang L. 等。Simulation-based route planning for pedestrian evacuation in metro stations: A case study。Automation in Construction，2016。
    \item Shi C. 等。Modeling and safety strategy of passenger evacuation in a metro station in China。Safety Science，2012。
    \item Xu X.-Y. 等。Analysis of subway station capacity with the use of queueing theory。Transportation Research Part C，2014。
\end{enumerate}

\end{document}
